\subsection{伪谱}

定义 9. --- 设 E 为复巴拿赫空间，u 为 E 上的闭部分算子。设 \(\varepsilon\) 为严格正的实数。我们把 u 的 \(\varepsilon\)-伪谱称为 u 的谱与如下的复数 \(\lambda\) 的集合之并：这些复数属于 u 的预解集且满足 \(\| \mathrm{R}(u,\lambda)\| >\varepsilon^{-1}\)。我们把这个集合记作 \(\mathrm{PSp}_{\varepsilon}(u)\)。

\pandocbounded{\includegraphics[keepaspectratio]{images/part002_14758b8b.jpg}}

有些作者把 \(\varepsilon\)-伪谱（对 \(u\)）定义为集合 \(\mathrm{T}_{\varepsilon}(u)\)，即 \(\mathrm{Sp}(u)\) 与如下的 \(\lambda \in \mathbf{C} - \mathrm{Sp}(u)\) 所成集合之并：它们满足 \(\| \mathrm{R}(u,\lambda)\| \geqslant \varepsilon^{-1}\)。\(\mathrm{PSp}_{\varepsilon}(u)\) 的闭包包含于 \(\mathrm{T}_{\varepsilon}(u)\)，但这一包含可以是严格的，即使 E 是希尔伯特空间亦然（参见 IV，第 348 页习题 18 及第 349 页习题 19）。

命题 18. --- 设 E 为复巴拿赫空间，u 为 E 上的闭部分算子。设 \(\varepsilon \in \mathbf{R}_{+}^{*}\)。集合 \(\mathrm{PSp}_{\varepsilon}(u)\) 是 C 的开子集。当 E 非零时它非空。

若 E 为零，则 \(\mathrm{PSp}_{\varepsilon}(u)\) 是空集。设 E 非零。由 IV，第 246 页命题 14，集合 \(\mathrm{PSp}_{\varepsilon}(u)\) 是 C 的开子集。

集合 \(\mathrm{PSp}_{\varepsilon}(u)\) 非空，如果 \(u\) 的谱非空。若 \(\mathrm{Sp}(u)\) 为空集，则刘维尔定理（VAR，R1，第 29 页，3.3.6）蕴含 \(\mathbf{C}\) 上由 \(\lambda \mapsto \mathrm{R}(u,\lambda)\) 定义的全纯函数不有界，故存在 \(\lambda\) 在 \(\mathbf{C}\) 中，使得 \(\| \mathrm{R}(u,\lambda)\| >\varepsilon^{-1}\)。

命题 19. --- 设 E 为复巴拿赫空间，u 为 E 上的闭部分算子。则我们有 \(\mathrm{PSp}_{\varepsilon}(u) \subset \mathrm{PSp}_{\delta}(u)\)（当 \(0 < \varepsilon < \delta\) 时），并且

\[
\bigcap_ {\varepsilon \in \mathbf {R} _ {+} ^ {*}} \mathrm{PSp} _ {\varepsilon} (u) = \mathrm{Sp} (u).
\]

第一个断言由定义立得。至于第二个断言，按定义，u 的谱包含于 \(\mathrm{PSp}_{\varepsilon}(u)\)，对一切 \(\varepsilon > 0\) 成立；而若 \(\lambda \notin \mathrm{Sp}(u)\)，则 \(\lambda \notin \mathrm{PSp}_{\varepsilon}(u)\)，只要 \(\varepsilon < \|R(u, \lambda)\|^{-1}\)。

命题 20. --- 设 E 为复巴拿赫空间，u 为 E 上的闭部分算子。设 \(\varepsilon\in\mathbf{R}_{+}^{*}\)。对每个 \(\lambda\in\mathbf{C}\)，下列条件等价：

\begin{enumerate}
\def\labelenumi{(\roman{enumi})}
\item
  有 \(\lambda \in \mathrm{PSp}_{\varepsilon}(u)\)；
\item
  或者 \(\lambda \in \mathrm{Sp}(u)\)，或者存在 \(x \in \mathrm{dom}(u)\) 满足 \(\| x \| = 1\) 且 \(\| (\lambda 1_{\mathrm{E}} - u)(x) \| < \varepsilon\)；
\item
  存在 \(v \in \mathcal{L}(\mathrm{E})\) 满足 \(\| v \| < \varepsilon\) 且 \(\lambda \in \mathrm{Sp}(u + v)\)。
\end{enumerate}

可以假设 E 非零。条件 (ii) 是定义的重新表述，因而也就是条件 (i)。

设条件 (i) 成立，我们来证明 (iii)。若 \(\lambda\) 属于 \(u\) 的谱，则在 (iii) 中可取 \(v = 0\)。

于是设 \(\lambda \notin \mathrm{Sp}(u)\)。由 \(\mathrm{PSp}_{\varepsilon}(u)\) 的定义，存在 \(y \in \mathrm{E}\) 满足 \(\| y \| = 1\) 且 \(\| \mathrm{R}(u, \lambda)y \| > \varepsilon^{-1}\)。令 \(x = \mathrm{R}(u, \lambda)y\)。我们有 \(x \neq 0\)。由哈恩–巴拿赫定理（EVT，II，第 67 页，推论 2），线性形式 \(\ell\)（在 \(\mathbf{C}x\) 上、满足 \(\ell(x) = 1\)）容许连续延拓 \(\ell_1 \in \mathrm{E}'\) 满足 \(\| \ell_1 \| = \| \ell \|\)；于是 \(\| \ell_1 \| = \| x \|^{-1} < \varepsilon\)。对每个 \(e \in \mathrm{E}\)，令 \(v(e) = \ell_1(e)y\)。我们有 \(v \in \mathcal{L}(\mathrm{E})\) 且 \(v(x) = y\)。我们得到 \((u + v)x = u(x) + y = \lambda x\)，故 \(\lambda \in \mathrm{Sp}(u + v)\)。又因 \(\| v \| = \| \ell_1 \| < \varepsilon\)，条件 (iii) 得到满足。

反之，设 \(v \in \mathcal{L}(\mathrm{E})\) 满足 \(\|v\| < \varepsilon\) 且 \(\lambda \in \mathrm{Sp}(u + v)\)。于是部分算子 \(\lambda1_{\mathrm{E}} - (u + v)\) 或者不是单射的，或者没有连续的逆；把 IV，第 243 页的引理 4 应用于单射部分算子 \(\lambda1_{E} - u\) 与 \(-v\)，便得 \(\|v\| \geqslant \|R(u, \lambda)\|^{-1}\)。由此推出条件 (iii) 蕴含 (i)。

推论。— 设 E 为复巴拿赫空间，u 为 E 上的闭部分算子。设 \(\varepsilon > 0\)。

\begin{enumerate}
\def\labelenumi{\alph{enumi})}
\item
  对每个 \(v \in \mathcal{L}(\mathrm{E})\)，我们有 \(\mathrm{PSp}_{\varepsilon}(u) \subset \mathrm{PSp}_{\varepsilon + \| v\|}(u + v)\)；
\item
  设 \(\delta > 0\)，\(D_{\delta}\) 为以 0 为心、以 \(\delta\) 为半径的 \(\mathbf{C}\) 中的开圆盘。我们有 \(\mathrm{PSp}_{\varepsilon}(u) + D_{\delta} \subset \mathrm{PSp}_{\varepsilon + \delta}(u)\)。
\end{enumerate}

设 \(\lambda \in \mathrm{PSp}_{\varepsilon}(u)\)。存在自同态 \(w \in \mathcal{L}(\mathrm{E})\) 满足 \(\|w\| < \varepsilon\) 且 \(\lambda \in \mathrm{Sp}(u+w)\)（命题 20）。由于 \(u+w = (u+v)+(w-v)\) 且 \(\|w-v\| < \varepsilon + \|v\|\)，我们有 \(\lambda \in \mathrm{PSp}_{\varepsilon+ \|v\|}(u+v)\)（同上所引）。

设 \(\mu \in \mathrm{D}_{\delta}\)；我们有 \(\lambda + \mu \in \mathrm{Sp}(u + (w + \mu 1_{\mathrm{E}}))\) 且 \(\| w + \mu 1_{\mathrm{E}} \| < \varepsilon + \delta\)，故 \(\lambda + \mu \in \mathrm{PSp}_{\varepsilon + \delta}(u)\)（同上所引）。

命题 21. --- 设 E 为复巴拿赫空间，u 为 E 上的闭部分算子。设 \(\varepsilon\in\mathbf{R}_{+}^{*}\)。\(\mathrm{PSp}_{\varepsilon}(u)\) 的每个有界连通分支都与 \(u\) 的谱相交。

设 U 是 \(\mathrm{PSp}_{\varepsilon}(u)\) 的一个不与 \(\mathrm{Sp}(u)\) 相交的连通分支。集合 U 在 \(\mathrm{PSp}_{\varepsilon}(u)\) 中既开又闭，因此它在 C 中的闭包 \(\overline{U}\) 满足 \(\overline{U} \cap \mathrm{PSp}_{\varepsilon}(u) = U\)。

由于 \(\mathrm{Sp}(u)\) 包含于 \(\mathrm{PSp}_{\varepsilon}(u)\) 而 U 与 \(\mathrm{Sp}(u)\) 不相交，这一等式表明集合 \(\overline{U}\) 与 \(\mathrm{Sp}(u)\) 不相交，从而包含于 u 的预解集。此外，它还蕴含集合 \(\overline{U}-U\) 与 \(\mathrm{PSp}_{\varepsilon}(u)\) 不相交。

于是我们有 \(\|\mathrm{R}(u,\lambda)\| \leqslant \varepsilon^{-1}\)，对一切 \(\lambda\)（属于 \(\overline{U} - U\)）；而 \(\|\mathrm{R}(u,\lambda)\| > \varepsilon^{-1}\)，对 \(\lambda \in U\)。若集合 U 有界，则其闭包 \(\overline{U}\) 是紧的，且存在 \(\lambda_{0} \in \overline{U}\) 使得 \(\|\mathrm{R}(u,\lambda)\| \leqslant \|\mathrm{R}(u,\lambda_{0})\|\) 对 \(\lambda \in \overline{U}\) 成立。前述结果表明 \(\lambda_{0} \in U\)，这与极大值原理（VAR，R1，第 29 页，3.3.7）矛盾，因为 u 的预解式在 u 的预解集中是全纯的（IV，第 246 页，命题 14）。

